\documentclass[11pt]{article}
\usepackage[margin=1in]{geometry}
\usepackage{amsmath,amssymb,amsthm}
\usepackage{booktabs}
\usepackage{listings}
\usepackage{xcolor}
\usepackage[numbers]{natbib}
\usepackage{hyperref}

\newtheorem{theorem}{Theorem}
\newtheorem{lemma}[theorem]{Lemma}
\theoremstyle{remark}
\newtheorem{remark}[theorem]{Remark}

\lstdefinelanguage{lean}{
  morekeywords={theorem,def,import,namespace,end,open,by,intro,exact,obtain,refine,have,simp,only,at,fun},
  sensitive=true, morecomment=[l]{--}, morecomment=[s]{/-}{-/}, morestring=[b]"}
\lstset{language=lean, basicstyle=\ttfamily\footnotesize, keywordstyle=\color{blue!60!black}\bfseries,
  commentstyle=\color{green!40!black}, columns=fullflexible, breaklines=true, frame=single,
  extendedchars=true, inputencoding=utf8,
  literate={≤}{{$\le$}}1 {≥}{{$\ge$}}1 {→}{{$\to$}}1 {↔}{{$\leftrightarrow$}}1 {ℝ}{{$\mathbb{R}$}}1
           {⟨}{{$\langle$}}1 {⟩}{{$\rangle$}}1 {∃}{{$\exists$}}1 {∀}{{$\forall$}}1 {∧}{{$\wedge$}}1
           {·}{{$\cdot$}}1 {−}{{$-$}}1}

\title{Re-certifying the Carbon Continuity Threshold:\\ A Lean Proof of the Two-Pool Regenerative Coupling Condition\\ under a Pinned Toolchain}
\author{Justin Hart\\ Viridis LLC}
\date{29 October 2026}

\begin{document}
\maketitle

\begin{abstract}
The live Viridis OS kernel \emph{Carbon Continuity} decides whether a two-pool carbon portfolio,
made of a living pool and a durable pool, can be kept from shrinking across disturbance-recovery cycles.
One cycle maps $(L,D)$ to $(aL+rD,\ pL+dD)$, with retention fractions $a,d\in[0,1)$ and transfer rates $r,p>0$.
The kernel's formal core states that a strictly positive portfolio that one cycle does not shrink in either pool
exists exactly when $(1-a)(1-d)\le rp$.
Until now this result rested on a certificate issued by a checker that has since been retired.
We restate the theorem as the sealed statement \texttt{P2.CarbonContinuity}
and give a short Lean~4 proof of it under Lean v4.28.0 with a pinned Mathlib.
The proof is ported from the archived canon release~\cite{HartCanon2026}.
We say exactly what the theorem settles, namely the threshold logic the kernel applies.
We also say what it does not settle: any coefficient value, the adequacy of the linear model,
and any ecological or carbon-credit outcome.
\end{abstract}

\section{Introduction}
Viridis OS runs a set of decision kernels. Each one is a small piece of software that turns ecological
parameters into a decision. The Carbon Continuity kernel
(runtime \texttt{viridis.carbon-continuity-threshold.v1}) answers one question: after a wildfire or a similar
disturbance, can a carbon portfolio be kept from shrinking cycle after cycle?
It answers by comparing a \emph{leakage product} $(1-a)(1-d)$ with a \emph{regenerative loop gain} $rp$.

The rule is only as trustworthy as the mathematics behind it. The canon release~\cite{HartCanon2026}
proved the rule in Lean, but the certificate came from tooling that can no longer be re-run.
This note re-proves the rule against a sealed restatement, so the kernel's mathematical assurance
can be re-checked with the current toolchain.
The question is narrow on purpose. We ask whether the threshold rule is a correct consequence of the stated model.
We do not ask whether the model or its coefficients describe any real forest.

\section{Background and related work}
The cycle map is the linear map $x\mapsto Mx$ with the nonnegative matrix
\[
M=\begin{pmatrix} a & r\\ p & d\end{pmatrix},\qquad x=\begin{pmatrix}L\\ D\end{pmatrix}.
\]
Since $r,p>0$, the matrix $M$ is irreducible.
By the Collatz--Wielandt characterisation in Perron--Frobenius theory~\cite{BermanPlemmons1994},
the spectral radius satisfies $\rho(M)\ge 1$ exactly when some $x>0$ has $Mx\ge x$.
In addition, $\rho(M)\ge 1$ holds, given $a,d<1$, exactly when $\det(I-M)=(1-a)(1-d)-rp\le 0$.
The theorem proved here is therefore a two-dimensional instance of classical theory.
We claim no novelty for the mathematics. The contribution is a re-checkable machine proof
tied to a live decision kernel.
The canon release~\cite{HartCanon2026} gives the ecological framing and the original Lean development,
including corollaries that the sealed statement does not include (boundary stationarity, strict growth above threshold).

\section{Model and statement}
The sealed module defines
\begin{align*}
\mathrm{stepLiving}(a,r,L,D) &= aL+rD, \qquad \mathrm{stepDurable}(d,p,L,D) = pL+dD,\\
\mathrm{Sustainable}(a,d,r,p) &:\iff \exists L,D,\ 0<L,\ 0<D,\ L\le aL+rD,\ D\le pL+dD.
\end{align*}

\begin{theorem}[\texttt{P2.CarbonContinuity}]
For all real $a,d,r,p$ with $0\le a<1$, $0\le d<1$, $r>0$ and $p>0$,
\[
\mathrm{Sustainable}(a,d,r,p)\iff (1-a)(1-d)\le rp .
\]
\end{theorem}

\section{Proof}
\begin{lemma}[Sufficiency]
If $a<1$, $r>0$ and $(1-a)(1-d)\le rp$, then $(L,D)=(r,1-a)$ is a witness.
\end{lemma}
\begin{proof}
Both components are positive. For the living pool, $ar+r(1-a)=r$, so the inequality holds with equality.
For the durable pool, $pr+d(1-a)-(1-a)=rp-(1-a)(1-d)\ge 0$.
\end{proof}

\begin{lemma}[Necessity]
If $a<1$, $d<1$ and $(L,D)$ is a witness, then $(1-a)(1-d)\le rp$.
\end{lemma}
\begin{proof}
The two inequalities rearrange to $(1-a)L\le rD$ and $(1-d)D\le pL$.
Both left-hand sides are positive, so we can multiply the inequalities:
$(1-a)(1-d)\,LD\le rp\,LD$. Cancelling $LD>0$ gives the threshold.
\end{proof}

\paragraph{Lean development.}
The proof file \texttt{proofs/P2.CarbonContinuity.lean} follows the canon proof almost line for line.
It is restated against the sealed definitions, and the heavier \texttt{nlinarith} calls are replaced
by \texttt{linarith} where the goal is linear after normalisation.
\begin{lstlisting}
theorem sufficient (a d r p : ℝ) (ha1 : a < 1) (hr : 0 < r)
    (h : (1 - a) * (1 - d) ≤ r * p) : Sustainable a d r p := by
  refine ⟨r, 1 - a, hr, by linarith, ?_, ?_⟩
  · simp only [stepLiving]; nlinarith
  · simp only [stepDurable]; nlinarith

theorem necessary (a d r p : ℝ) (ha1 : a < 1) (hd1 : d < 1)
    (hs : Sustainable a d r p) : (1 - a) * (1 - d) ≤ r * p := by
  obtain ⟨L, D, hL, hD, hliv, hdur⟩ := hs
  simp only [stepLiving] at hliv
  simp only [stepDurable] at hdur
  have h1 : (1 - a) * L ≤ r * D := by linarith
  have h2 : (1 - d) * D ≤ p * L := by linarith
  have h3 : ((1 - a) * L) * ((1 - d) * D) ≤ (r * D) * (p * L) :=
    mul_le_mul h1 h2 (by nlinarith) (by nlinarith)
  have h4 : ((1 - a) * (1 - d)) * (L * D) ≤ (r * p) * (L * D) := by linarith
  exact le_of_mul_le_mul_right h4 (mul_pos hL hD)

theorem main : ViridisStatements.P2.CarbonContinuity.Statement := by
  intro a d r p _ ha1 _ hd1 hr _
  exact ⟨necessary a d r p ha1 hd1, sufficient a d r p ha1 hr⟩
\end{lstlisting}

\section{Results}
We give a Lean proof of the sealed statement \texttt{P2.CarbonContinuity}.
The local pre-check tool reports that the file compiles under the pinned toolchain,
matches the sealed statement, and depends only on \texttt{propext}, \texttt{Classical.choice} and \texttt{Quot.sound}.
The labels are awarded by the independent verifier, whose certification page comes before this paper.

\begin{remark}
We observe that the helper lemmas use only $a<1$, $d<1$ and $r>0$.
The hypotheses $0\le a$, $0\le d$ and $p>0$ are part of the kernel's parameter domain,
but the proof does not need them.
\end{remark}

\section{What the theorem settles and what it does not}
\begin{itemize}
\item \textbf{Settled (subject to verification):} inside the stated linear two-pool model, the kernel's rule
``sustainable iff $(1-a)(1-d)\le rp$'' is a correct decision procedure, including at the boundary.
\item \textbf{Not settled:} the values of $a,d,r,p$ for any site. These come from field evidence and stay ASSUMED
until a TESTED lane or a human review covers them.
\item \textbf{Not settled:} whether a linear, time-invariant two-pool cycle model describes post-fire carbon
dynamics. The theorem also says nothing about uncertainty in the coefficients, about stochastic disturbance,
or about any carbon-credit or treatment outcome.
\item \textbf{Not settled:} the statement is about one cycle. Persistence over many cycles follows informally
because $M$ is monotone on nonnegative vectors, so $Mx\ge x$ implies $M^{k}x\ge x$. That fact is not part of the
sealed statement.
\item \textbf{Not settled:} that the TypeScript implementation (\texttt{src/lib/carbonContinuity.ts})
implements this rule faithfully. Code conformance has not been checked here.
\end{itemize}

\section{Limitations}
The result is mathematically elementary. Its value lies in being re-checkable and tied to a deployed kernel,
not in its depth. The canon corollaries (boundary stationarity, strict growth, nonvacuity) were not re-proved
as sealed statements. The sealed module includes concrete sanity examples that pin the meaning of the definitions.

\section{Reproducibility}
Toolchain \texttt{leanprover/lean4:v4.28.0}, Mathlib commit \texttt{8f9d9cff6bd728b17a24e163c9402775d9e6a365}.
Sealed statement source sha256 \texttt{a5a6e8c63205b865da0df0c302f1926b48d1874ce1310053da6c1ae3104ffded}.
Canon source sha256 \texttt{899d3bb1929bac8d1ab6b760caffac428de115489e6fc7c1b2752e0769a77fc1}
(canon repository commit \texttt{80d32904d751b45f415e23f872510cc2eddca3ff}).
To check the proof, run \texttt{viridis-lean-check proofs/P2.CarbonContinuity.lean P2.CarbonContinuity}.
No automated prover jobs and no external data were used.

\section{Changelog}
\begin{itemize}
\item 2026-10-29 (P2-T001): First version. Ported the canon threshold proof to the sealed statement
\texttt{P2.CarbonContinuity}, and wrote the scope statement.
\end{itemize}

\bibliographystyle{plainnat}
\bibliography{refs}
\end{document}
